\answer
\begin{enumerate}

%% 3章 問1（3.4節）
\item[3.4節 問1]
(a) 原点付近の直線の傾きの逆数がオン抵抗だから
$R_{\mathrm{on}} = v_{DS}/i_D = 0.5/6 \approx 0.083\,\Omega = 83\,\mathrm{m\Omega}$。
ゲート電圧が高いほど反転層の電子が多くオン抵抗は小さいので，
スイッチとして使うときはゲートをしきい値よりずっと高く（$10$〜$15$\,V程度）駆動する。
(b) ゲートは酸化膜（絶縁体）の上にあるので，
定常状態ではゲートに電流が流れず，駆動電力はほぼゼロである。
一方，ゲート電極と半導体は酸化膜をはさんだキャパシタ（ゲート容量）をなすから，
オン・オフを切り替える瞬間には，
このキャパシタを充電・放電するための電流がゲートに流れる。
したがって駆動電力はスイッチング周波数に比例して増える。

%% 3章 問2（3.5節）
\item[3.5節 問2]
IGBTの主電流は，裏面のp形コレクタ層とn$^-$ドリフト層の間のpn接合を
順方向に通り抜ける。
順バイアスのpn接合には約$0.7$\,Vの順方向電圧が必ず立つから（\ref{sec3.2}節），
IGBTのオン電圧にはこの分の下駄が乗る。
MOSFETでは電子がソースから反転層を通ってドレインまで
pn接合の障壁を一度も越えずに流れるため，この下駄がない。

%% 3章 問3（3.7節）
\item[3.7節 問3]
アバランシェ降伏は，電場で加速されたキャリアが価電子帯の電子を
伝導帯へたたき上げる衝突電離の連鎖で起こる。
1回の衝突電離で電子正孔対を作るには，
キャリアがバンドギャップ$E_g$以上の運動エネルギーをもつ必要がある。
$E_g$が大きい材料では，衝突と衝突の間の短い距離で
それだけのエネルギーをキャリアに与えるために，
より強い電場が必要になる。
したがって$E_g$が大きいほど絶縁破壊電界$\mathcal{E}_c$は大きい。

%% 3章 問4（3.9節）
\item[3.9節 問4]
式\ref{eq3.5}より
\begin{equation*}
W_{\mathrm{GaN}} = \frac{2\times1200}{3.3\times10^8}
\approx 7.3\times10^{-6}\,\mathrm{m} = 7.3\,\mathrm{\mu m}
\end{equation*}
シリコンの$80\,\mathrm{\mu m}$（例題\ref{rei3.2}）の約11分の1である。
絶縁破壊電界が11倍大きい分だけ，同じ耐圧を受け持つ層は11分の1に薄くできる。

\end{enumerate}
